\chapter{Biological inspiration and its limits}
\label{ch:biology-limits}

\section*{A useful analogy must be allowed to fail}
Chapter~\ref{ch:what-adds} showed that selective execution is not uniquely
genomic. Does that make biology irrelevant to computational design? No. It
means we need a more careful method for borrowing an idea. We should say what
we observed, what abstraction we extracted, what we omitted and which
experiment could show that the abstraction was unhelpful.

Three activities must remain distinct. An \emph{analogy} highlights a possible
similarity. An \emph{abstraction} specifies a smaller mathematical object with
explicit rules. A \emph{biological model} makes testable claims about a biological
system. An abstraction inspired by regulation can be a useful computer program
without being an adequate biological model. Conversely, biological realism
does not automatically make a computational design efficient.

This chapter uses one biological distinction and one original recurrence.
The distinction separates sequence transfer from regulatory control. The
recurrence separates an instantaneous response from persistent state. Neither
establishes that Evolutor has a genome, developmental program or learning
mechanism.

\section{Two meanings that one arrow can conceal}
\index{regulatory control}\index{sequence transfer}
For a protein-coding pathway, the familiar DNA-to-RNA-to-protein sketch concerns
the relationship between polymer sequences. Crick's 1970 discussion explicitly
distinguished such \gls{sequence-transfer} from control over the rates of processes;
reverse transcription was not excluded by his formulation
\citep{crick1970}. A diagram that merges these arrow meanings can create a
false contradiction where none exists.

Proteins can regulate transcription through interactions with DNA control
sequences. The classical distinction between a linked control site and a
diffusible regulatory factor is useful here \citep{cooper2000transcription}.
A regulatory protein changing transcription is not its amino-acid sequence
being translated backward into DNA. This is the narrow point of
Figure~\ref{fig:information-control}; it is not an exhaustive map of molecular
biology.

\begin{figure}[p]
\centering
\includegraphics[width=\textwidth]{../build/figures/evo-g07.pdf}
\caption{EVO-G07. Sequence transfer and regulatory control are different
relations. The schematic uses a protein-coding pathway and a generic
transcriptional repressor, not a complete operon or a universal feedback loop.
Solid transfer arrows and the dashed inhibitory line must not be interpreted
interchangeably \citep{crick1970,cooper2000transcription}.}
\label{fig:information-control}
\end{figure}

For software, this distinction suggests a question rather than an answer:
does a component supply content, alter a control decision, or both? A module
returning a value, a gate selecting a module, and an updater changing a gate
are different mechanisms. Giving every connecting arrow the label
``information'' makes the architecture harder to audit. Typed relations expose
what would need to be implemented and tested.

\section{State is not the same as an enabled operation}
An instantaneous gate answers a question about the current input. A stateful
response also depends on what happened before. To isolate this difference,
define a nonnegative state $x_t$, a binary signal $u_t\in\{0,1\}$, a production
parameter $p\geq0$ and a loss fraction $d\in[0,1]$. Our discrete update is
\begin{equation}
 x_{t+1}=(1-d)x_t+p u_t.
 \label{eq:leaky-state}
\end{equation}
The comparison gate returns $y_t=p u_t$ with no retained state. We compare
$y_t$ with $x_{t+1}$ after processing the same input $u_t$, not with the state
before that input. The indexing is part of the experiment.

This is an ordinary \gls{leaky-integrator} written as an original teaching example.\index{leaky integrator}
Its units are arbitrary state units per discrete update. It is not fitted to
RNA, protein or any measured cell population. The word ``production'' is a
mnemonic for the additive term, not evidence that the equation captures an
expression mechanism. In particular, $d$ is a fraction removed per update,
not a continuous-time decay constant.

Let $p=2$, $d=1/2$, $x_0=0$ and use the signals $(1,1,0,0,0,1)$.
Table~\ref{tab:state} is generated using exact rational arithmetic.
After two active steps, the signal becomes zero. The instantaneous output
immediately becomes zero, but the state remains positive and then decays.
When the signal returns, the retained remainder contributes to the new state.

\begin{table}[htbp]
\centering
% Generated by scripts/chapter02_artifacts.py; do not edit.
\begin{tabular}{rrrrr}
\toprule $t$ & $u_t$ & $x_t$ & $x_{t+1}$ & $y_t$ \\
\midrule
0 & 1 & $0$ & $2$ & $2$ \\
1 & 1 & $2$ & $3$ & $2$ \\
2 & 0 & $3$ & $3/2$ & $0$ \\
3 & 0 & $3/2$ & $3/4$ & $0$ \\
4 & 0 & $3/4$ & $3/8$ & $0$ \\
5 & 1 & $3/8$ & $35/16$ & $2$ \\
\bottomrule\end{tabular}

\caption{An exact synthetic pulse response in arbitrary state units.
The table reports the state before and after each discrete input, alongside
an instantaneous gate. No biological measurement or fitted rate is used.}
\label{tab:state}
\end{table}

At $t=2$ and $t=4$, the current input is zero in both cases, but the resulting
states differ. No function of that binary current input alone can reproduce
both outputs. This establishes a distinction between the declared stateless
and stateful models; it does not establish a uniquely biological or genomic
property. A conventional recurrent program handles it directly.

\section{A closed form makes the memory explicit}
Put $a=1-d$. Repeated substitution in Equation~\ref{eq:leaky-state} yields
\begin{equation}
 x_t=a^t x_0+p\sum_{j=0}^{t-1}a^{t-1-j}u_j.
 \label{eq:state-convolution}
\end{equation}
For $t=0$ the sum is empty. A signal's contribution is weighted by the number
of updates since it occurred. This is the exact reason history matters here,
not an inference from the biological vocabulary.

\begin{proposition}[History weighting and a bound]
Equation~\ref{eq:state-convolution} solves the recurrence. If $0<d\leq1$,
$x_0\geq0$ and signals are binary, then
$0\leq x_t\leq\max\{x_0,p/d\}$ for every $t$.
\end{proposition}
\begin{proof}
The expression gives $x_0$ at $t=0$. Multiplying its value at $t$ by $a$ and
adding $p u_t$ produces the stated expression at $t+1$, proving the formula
by induction. For the bound, let $M=\max\{x_0,p/d\}$. If $0\leq x_t\leq M$,
then $0\leq a x_t+p u_t\leq(1-d)M+p\leq M$, because $dM\geq p$.
The initial condition starts the induction.
\end{proof}

For a run of zero inputs, the state after $k$ additional updates is $a^k$
times its initial value for that run. With $d=1/2$, it halves at each update.
With $d=1$, the previous state is completely removed and the model reduces to
the instantaneous gate. With $d=0$, nothing is removed: repeated active inputs
accumulate without the finite bound above. Keeping these boundary cases in
view prevents a proof valid for $d>0$ from being extended to $d=0$ by accident.

For $0<d<1$, arbitrarily old inputs have nonzero weights in exact arithmetic,
although their influence decreases. Thus ``memory'' here does not mean a
finite window in which all older inputs are exactly forgotten. The state is
finite-dimensional; its dependency on history is not finitely truncated.

\section{What would make this a biological model?}
To interpret $x_t$ as a measured biological quantity, we would need to choose
that quantity and its units, define an observation process, select a time
interval and justify the update law. A measurement of total signal would not
automatically identify the separate contributions represented by production
and loss. We would need experiments that distinguish competing explanations.

Our synthetic sequence cannot do that job. Its parameters were chosen for
readable fractions, not estimated from data. Even a visually plausible curve
would not validate the mechanism. Different hidden processes can agree over
a short observation window. The next useful step would be a discriminating
experiment, not a more elaborate biological name for the same recurrence.

The same caution applies in machine learning. If retained state improves a
sequence task, compare with an ordinary recurrent baseline under the same
information and resource constraints. If state is reset between evaluation
examples, report that policy. If it persists, ensure that it cannot carry
information from a held-out answer into later predictions. State is both a
modeling resource and a possible channel for leakage.

\section{Separate changes that happen at different levels}
For the computational proposals later in this book, we distinguish three
update categories. A \emph{state update} changes a run's carried values, as
Equation~\ref{eq:leaky-state} does. A \emph{parameter update} changes a rule's
coefficients or selection parameters. A \emph{structural update} changes the
available operations or their connections. These are design categories, not
a claim that biology divides neatly into three independent clocks.

Changing $x_t$ with fixed $p$ and $d$ is not learning those parameters. Adjusting
$p$ by an optimizer does not add a new operation to a library. Adding an
operation is not inheritance unless the system specifies what is transmitted
between which entities. This vocabulary blocks several attractive but invalid
shortcuts: dynamics do not by themselves establish development; parameter
fitting does not by itself establish evolution; persistence does not by itself
establish a genome.

A candidate mechanism should therefore name the object that changes and the
evidence that drives the change. If a proposed system updates all three levels,
each needs its own state transition, test and resource account. Otherwise a
reported benefit cannot be attributed to the mechanism the author highlights.
This is the same principle as Chapter 1's routing-cost account, now applied to
adaptation rather than invocation.

\section{Reproduce the calculation and inspect its limits}
\begin{Verbatim}[fontsize=\small]
from minievolutor.leaky_state import chapter_example
for step in chapter_example():
    print(step.time, step.signal, step.before, step.after)
\end{Verbatim}
The code accepts integer or rational parameters, rejects floating-point
parameters rather than silently treating decimal approximations as exact,
and validates the domain. Tests compare recurrence outputs with the explicit
convolution for every binary input sequence of lengths zero through seven,
at three loss settings: 765 input/setting cases. They also check the zero-input
decay, nonnegativity and the two boundary loss fractions.

The graph source identifies which arrows describe biology and which relationships
belong only to the computational example. There is no learned gate, stochastic
reaction, optimizer, mutation operator or external action in this code. Those
omissions are visible scope decisions, not capabilities waiting to be inferred.

\section{Exercises and solution notes}
\begin{enumerate}
\item Starting from $x_0=8$ with $d=1/2$ and three zero inputs, compute the final
state. Why does the production parameter not affect this answer?
\item With $d=0$, $p=2$ and zero initial state, compute the response to
$(1,0,1)$. Which bound in the proposition is unavailable?
\item Construct two input histories ending in the same current signal but
yielding different states. State which model class this distinguishes.
\item Explain why a repressor-to-DNA control line is not a reverse-translation
arrow. What information would an unlabeled arrow conceal?
\item An author calls the change from $x_t$ to $x_{t+1}$ ``genome evolution.''
What object would need to change, and what inheritance/variation rules are
missing from our example?
\item A new architecture beats an instantaneous gate on the pulse example.
Name a stronger ordinary baseline and a test that would isolate the claimed
contribution beyond state retention.
\end{enumerate}

\paragraph{Solution notes.}
(1) One; all production terms are multiplied by zero. (2) The post-update
states are 2, 2, 4; $p/d$ is undefined. (3) Histories $(1,0)$ and $(0,0)$ give
post-update states 1 and 0 under the default parameters. They distinguish this
stateful recurrence from a deterministic function of the current signal alone.
(4) Control affects a process rather than encoding the protein's residue
sequence in DNA; the relation type is essential. (5) The example has no genome,
inherited structure or variation/selection rule; only a scalar state changes.
(6) Compare with the same leaky integrator and ablate the proposed additional
mechanism at matched budgets and state-access rules.
